
\section{Kinect}

\begin{frame}[b]
	\begin{figure}
		\includegraphics[width=0.8\textwidth]{nav_vision}
	\end{figure}
	\vspace*{0.7cm}
\end{frame}

\begin{frame}[t]{Kinect}
	\begin{block}{Koordinatensystem}
		\begin{columns}[t]
			\begin{column}{0.4\textwidth}
				\begin{itemize}
					\item Koordinatenursprung im Tiefensensor
					\item Treiber liefert Angaben in Millimetern
					\item Distanzen werden auf Kameraebene projeziert wiedergegeben
				\end{itemize}
			\end{column}
			\begin{column}{0.4\textwidth}
				\begin{figure}
					\includegraphics[width=0.8\textwidth]{vision_coordsys}
				\end{figure}
			\end{column}
		\end{columns}
	\end{block}
\end{frame}

\begin{frame}[t]{Kinect}
	\begin{block}{API}
		\begin{itemize}
			\item Auswahl zwischen Microsoft SDK, OpenKinect, Freenect und
				OpenNI
			\item Entscheidung für OpenNI gemäß \cite{andersen2012kdsefcva}
			\item Umwandlung der Bilddaten in OpenCV Datenstrukturen
		\end{itemize}
	\end{block}
\end{frame}

\begin{frame}[t]{Kinect}
	\begin{block}{Probleme und Lösungen}
		\begin{itemize}
			\item Umwandlung der Bilddaten teuer \arrow ?
			\item Tiefenwerte naher Objekte werden nicht erkannt \arrow
				Mindestabstand von 50cm (Linux) bzw. 80cm (Windows)
		\end{itemize}
	\end{block}
\end{frame}

